\chapter*{Conclusion}
\addcontentsline{toc}{chapter}{Conclusion}

\section*{Summary of Achievements}
The development of the Lectura platform demonstrates the immense potential of integrating state-of-the-art Artificial Intelligence into modern educational workflows. As information overload continues to emerge as a critical challenge for university students, tools that facilitate active learning and knowledge extraction become paramount. 

Lectura successfully achieved its primary objective of providing an integrated, multimodal AI study assistant. By supporting diverse inputs such as YouTube lecture videos, PDF documents, and plain text, the system removes the initial friction of content ingestion. The downstream transformation pipeline then effectively applies Large Language Models (specifically Google Gemini 3 Flash) to generate pedagogically grounded study materials, including varied text summaries, multiple-choice quizzes, spaced repetition flashcard decks, and context-aware conversational chat responses.

From a software engineering perspective, the project validated the architectural viability of a unified Go-based asynchronous backend communicating with a highly reactive React/TypeScript frontend. The system manages long-running AI inference tasks robustly using a Redis-backed worker pool, pushing real-time status updates via WebSockets. Combining PostgreSQL for reliable transactional data and JWT rotation with Google OAuth for security, the solution meets modern web application standards.

\section*{Fulfillment of Objectives}
Revisiting the core technical and functional objectives:
\begin{itemize}[label=\textendash]
    \item \textit{Multimodal Ingestion}: Fully implemented. Text extraction from documents and YouTube API transcript retrieval function smoothly.
    \item \textit{Generative Transformation}: The platform successfully utilizes structured prompts to yield strict JSON outputs mapping sequentially to frontend components (Summaries, Quizzes, Flashcards).
    \item \textit{Interactive Testing}: Users are able to take self-assessments directly on the platform and log their success rates within the PostgreSQL database.
    \item \textit{PDF Export}: Material generated on the platform is rendered into structured portable documents for offline review.
\end{itemize}

\section*{Future Work}
While Lectura provides a comprehensive minimum viable product, several avenues for future enhancement have been identified:
\begin{itemize}[label=\textendash]
    \item \textit{Multimodal Ingestion}: Fully implemented. Text extraction from documents and YouTube API transcript retrieval function smoothly.
    \item \textit{Generative Transformation}: The platform successfully utilizes structured prompts to yield strict JSON outputs mapping sequentially to frontend components (Summaries, Quizzes, Flashcards).
    \item \textit{Interactive Testing}: Users are able to take self-assessments directly on the platform and log their success rates within the PostgreSQL database.
    \item \textit{PDF Export}: Material generated on the platform is rendered into structured portable documents for offline review.
\end{itemize}

Ultimately, Lectura establishes a strong foundation for the next generation of EdTech platforms, proving that when AI is utilized not as a general-purpose oracle, but as a specialized structured transformation engine, it can significantly enhance academic productivity.
